\documentclass[12pt,letterpaper,twoside]{article}
\usepackage[utf8]{inputenc}
\usepackage[backend=biber, style=ieee]{biblatex}
\usepackage{amsmath}
\usepackage{listings}
\usepackage{graphicx}
\usepackage{amssymb}
\usepackage{array}
\usepackage[hidelinks]{hyperref}
\usepackage{multirow}
\usepackage{relsize}
\usepackage{placeins}
\usepackage{enumitem}
\usepackage{ragged2e}
\usepackage{csquotes}
\usepackage[spanish, english]{babel} % Cargar ambos idiomas
\usepackage{siunitx}
\usepackage{textcomp}

\usepackage[T1]{fontenc} 
\addbibresource{Referencias.bib}
\renewcommand{\theequation}{\thesubsubsection.\arabic{equation}} 
\usepackage[margin=2.5cm]{geometry}
\usepackage{multirow}
\usepackage{array}
\usepackage{xcolor}
\newenvironment{Encuesta}{\begin{figure}[h]\centering}{\end{figure}}

\renewcommand{\theequation}{\arabic{equation}}
\setcounter{equation}{0}

\begin{document}

\begin{center}
    \text{\large UNIVERSIDAD TECNOLÓGICA DE LA MIXTECA} \\
    \text{INSTITUTO DE ELECTRÓNICA Y MECATRÓNICA} \\[0.5cm]
    \text{\large Primer Examen Parcial de Control No Lineal} \\
    \text{MEOSIA} \\[0.5cm]
    \text{30 de septiembre de 2026} \\[0.5cm]
\end{center} \par

\textbf{Nombre:} Yahir Fernando García García     \hspace{5cm}    \textbf{Grupo:} MEOSIA3\par\vspace{0.1cm}
\textbf{Instrucciones:} La solución del examen debe incluir las demostraciones solicitadas escritas a mano, junto con los códigos de MATLAB y las simulaciones que validen los resultados. Se verificará la originalidad del trabajo
presentado.

\small
\vspace{0.5cm}
Considere el siguiente modelo de un PVTOL

\begin{equation*}
    \begin{aligned}
        \ddot{x}&=-\frac{U}{m}sin(\theta)\\
        \ddot{z}&=\frac{U}{m}cos(\theta)-g\\
        \ddot{\theta}&=\frac{1}{J}\tau
    \end{aligned}
\end{equation*}

Donde \(\theta\) es el ángulo de actitud. \({x}\) y \({z}\) representan las posiciones traslacionales, \(U\) es el empuje total generado por los motores, \(\tau\) es el par aplicado para controlar la actitud, \(m\) es la masa del vehiculo, \(J\) es el momento de inercia y \(g\) es la gravedad. Se usarán como parámetros \(m = 1\), \(J = 1\).

\begin{enumerate}
    \item Mediante linealización exacta por realimentación, realice el desacoplamiento de las dinámicas mediante entradas virtuales y obtenga el modelo que permita establecer un lazo de control externo para las posiciones (1 punto).
\end{enumerate}

Se define el vector de estados como el siguiente

\begin{equation*}
    X=\begin{bmatrix}
        x_1\\x_2\\x_3\\x_4\\x_5\\x_6
    \end{bmatrix}=\begin{bmatrix}
        x\\\dot{x}\\z\\\dot{z}\\\theta\\\dot{\theta}
    \end{bmatrix}
\end{equation*}

El sistema se escribe como

\begin{equation*}
    \dot{X}=f(X)+g_1(X)U+g_2(X)\tau
\end{equation*}

con:

\begin{equation*}
    f(X)=\begin{bmatrix}
        x_2\\0\\x_4\\-g\\x_6\\0
    \end{bmatrix},\quad g_1(X)=\begin{bmatrix}
        0\\-sin(x_5)\\0\\cos(x_5)\\0\\0
    \end{bmatrix},\quad g_2(X)=\begin{bmatrix}
        0\\0\\0\\0\\0\\1
    \end{bmatrix}
\end{equation*}

Se derivan las salidas hasta que aparezcan las entradas \(U\) y \(\tau\). Para la salida \(y_1=x_1\) se tiene que:

\begin{equation*}
    \begin{aligned}
        \dot{y}_1&=x_2\\
        \ddot{y}_1&=\dot{x}_2=-U sin(x_5)\\
    \end{aligned}
\end{equation*}

Dado que aparece U en la segunda derivada, se tiene que el orden relativo de la salida \(y_1\) es 2. Para la salida \(y_2=x_3\) se tiene que:

\begin{equation*}
    \begin{aligned}
        \dot{y}_2&=x_4\\
        \ddot{y}_2&=\dot{x}_4=U cos(x_5)-g\\
    \end{aligned}
\end{equation*}

Dado que aparece U en la segunda derivada, se tiene que el orden relativo de la salida \(y_2\) es 2. La suma de los grados relativos es 4, mientras que el orden del sistema es 6, por lo que se tiene un sistema no completamente linealizable. Las derivadas de las salidas se escriben como:

\begin{equation*}
    \begin{bmatrix}
        \ddot{y}_1\\\ddot{y}_2
    \end{bmatrix}=\begin{bmatrix}
        -sin(x_5) & 0\\cos(x_5) & 0
    \end{bmatrix}\begin{bmatrix}
        U\\\tau
    \end{bmatrix}+\begin{bmatrix}
        0\\-g
    \end{bmatrix}
\end{equation*}

por lo que la matriz de desacoplamiento seria:

\begin{equation*}
    A(X)=\begin{bmatrix}
        -sin(x_5) & 0\\cos(x_5) & 0
    \end{bmatrix}
\end{equation*}

Para desacoplar el sistemas, se definen las siguientes entradas virtuales:

\begin{equation*}
    v_x\to \ddot{x}, \quad v_z\to \ddot{z}
\end{equation*}

Sustituyendo en las expresiones del modelo tenemos que:

\begin{equation*}
    \begin{aligned}
        v_x&=-U sin(\theta)\\
        v_z&=U cos(\theta)-g
    \end{aligned}
\end{equation*}

sumando ambos terminos y despejando \(U\) se obtiene que:

\begin{equation*}
    \begin{aligned}
        v_x^2+(v_z+g)^2&=U^2 sin^2(\theta)+U^2 cos^2(\theta)\\
        U&=\sqrt{v_x^2+(v_z+g)^2}\\
    \end{aligned}
\end{equation*}

Dividiendo la primera ecuacion entre la segunda

\begin{equation*}
    \begin{aligned}
        \frac{v_x}{v_z+g}&=\frac{-U sin(\theta)}{U cos(\theta)}=-tan(\theta)\\
        \theta_d&=atan\left(-\frac{v_x}{v_z+g}\right)\\
    \end{aligned}
\end{equation*}

La dinámica de la actitud es:

\begin{equation*}
    \ddot{\theta}=\frac{1}{J}\tau
\end{equation*}

Entonces se define el error de actitud como \(e_\theta=\theta-\theta_d\), derivando dos veces se obtiene que:

\begin{equation*}
    \begin{aligned}
        \dot{e}_\theta&=\dot{\theta}-\dot{\theta}_d\\
        \ddot{e}_\theta&=\ddot{\theta}-\ddot{\theta}_d\\
        \ddot{e}_\theta&=\frac{1}{J}\tau-\ddot{\theta}_d\\
    \end{aligned}
\end{equation*}

Entonces se elige la ley de control para la actitud como:

\begin{equation*}
    \tau=J(\ddot{\theta}_d-k_1\dot{e}_\theta-k_2 e_\theta)
\end{equation*}

Sustituyendo \(\theta=\theta_d\) en las ecuaciones de translación, se obtiene el modelo desacoplado para el lazo externo:

\begin{equation*}
    \begin{aligned}
        \ddot{x}&=-U sin(\theta)\\
        \ddot{z}&=U cos(\theta)-g\\
    \end{aligned}
\end{equation*}

Por lo tanto, el modelo desacoplado es:

\begin{equation*}
    \begin{aligned}
        \ddot{x}&=v_x\\
        \ddot{z}&=v_z\\
    \end{aligned}
\end{equation*}

Con ello, el lazo externo queda definido por las siguientes ecuaciones:

\begin{equation*}
    \begin{aligned}
        U=\sqrt{v_x^2+(v_z+g)^2}\\
        \theta_d=atan\left(-\frac{v_x}{v_z+g}\right)\\
        \tau=J(\ddot{\theta}_d-k_1\dot{e}_\theta-k_2 e_\theta)\\
    \end{aligned}
\end{equation*}














\begin{enumerate}
    \item Para el sistema obtenido mediante el procedimiento anterior, diseñe e implemente los siguientes controladores para
el seguimiento de trayectorias en x y z:
    \begin{itemize}
        \item \textbf{Realimentación de estados usando LMI} (1 punto)\\
        
        Del punto 1, el desacoplamiento por linealización es:

        \begin{equation*}
            \ddot{x}=v_x,\quad \ddot{z}=v_z \quad \ddot{\theta}=\tau
        \end{equation*}
        
        Cada uno es un doble integrador. En espacio de estados se puede escribir como:

        \begin{equation*}
            \dot{x}_p=Ax_p+Bu_p
        \end{equation*}

        donde:

        \begin{equation*}
            A=\begin{bmatrix}
                0 & 1\\
                0 & 0
            \end{bmatrix},\quad B=\begin{bmatrix}
                0\\
                1
            \end{bmatrix}
        \end{equation*}

        Se propone la siguiente ley de control:

        \begin{equation*}
            u_p=-Kx_p
        \end{equation*}

        donde \(K=[k1\quad k2]\). Cerrando el lazo, se tiene que:

        \begin{equation*}
            \dot{x}_p=(A-BK)x_p
        \end{equation*}

        Para esto, se propone la siguiente función candidata de Lyapunov:

        \begin{equation*}
            V(x_p)=x_p^TPx_p
        \end{equation*}

        Derivando y sustituyendo, se obtiene que:

        \begin{equation*}
            \begin{aligned}
                \dot{V}(x_p)=&\dot{x}_p^TPx_p+x_p^TP\dot{x}_p\\
                =&x_p^T(A-BK)^TPx_p+x_p^TP(A-BK)x_p\\
                =&x_p^T[(A-BK)^TP+P(A-BK)]x_p
            \end{aligned}
        \end{equation*}

        Para que el sistema sea estable, se requiere que \(\dot{V}(x_p)<0\), por lo que se tiene que:

        \begin{equation*}
            (A-BK)^TP+P(A-BK)<0
        \end{equation*}

        Dado que hay productos entre \(P\) y \(K\), se realiza el cambio de variable:

        \begin{equation*}
            \begin{aligned}
                F&=KP^{-1}\\
                K&=FP
            \end{aligned}
        \end{equation*}

        Sustituyendo en la ecuación de estabilidad, se obtiene que:

        \begin{equation*}
            (A-BFP)^TP+P(A-BFP)<0
        \end{equation*}   

        sustituyendo

        \begin{equation*}
            \begin{aligned}
                A-BK=A-BFP\\
            \end{aligned}
        \end{equation*}

        Expandiendo

        \begin{equation*}
            \begin{aligned}
                (A-BFP)^TP+P(A-BFP)=A^TP+PA-P^TF^TB^TP-PBFP\\
            \end{aligned}
        \end{equation*}

        Dado que \(P\) es simétrica, se tiene que \(P^T=P\), por lo que:

        \begin{equation*}
            A^TP+PA-P^TF^TB^TP-PBFP=A^TP+PA-P(F^TB^T+BF)P\\
        \end{equation*}

        Haciendo el siguiente cambio de variable:

        \begin{equation*}
            \gamma=P^{-1},\quad F=K\gamma
        \end{equation*}

        Multiplicando por \(\gamma\) a la izquierda y derecha, se obtiene que:

        \begin{equation*}
            \begin{aligned}
                \gamma[A^TP+PA-P(F^TB^T+BF)P]\gamma<0\\
                \gamma A^TP\gamma+\gamma PA\gamma-\gamma P(F^TB^T+BF)P\gamma<0\\
            \end{aligned}
        \end{equation*}

        Dado que \(\gamma P = I\) y \(P\gamma = I\), se tiene que:

        \begin{equation*}
            \gamma A^T + A\gamma - (F^TB^T + BF) < 0
        \end{equation*}

        Dado que es necesario asegurarse que el contrador de \(\tau\) sea más rapido, se agrega un término de amortiguamiento, por lo que la ecuación final es:

        \begin{equation*}
            \gamma A^T + A\gamma - (F^TB^T + BF) + \alpha\gamma < 0
        \end{equation*}








        \item \textbf{Modos deslizantes} (2 puntos)\\

        Cada eje del modelo desacoplado se escribe como un sistema de segundo orden de la forma:

        \begin{equation*}
            \begin{aligned}
                \dot{X}_1 &= X_2 \\
                \dot{X}_2 &= f(X,t)+v
            \end{aligned}
        \end{equation*}

        Se desea que la salida $y=X_1$ siga una señal de referencia $r(t)$. Se define el error de seguimiento como:

        \begin{equation*}
            e = r-y = r-X_1
        \end{equation*}

        de modo que:

        \begin{equation*}
            \dot{e}=\dot{r}-X_2
        \end{equation*}

        Se propone la superficie de deslizamiento como una ecuación diferencial estable de primer orden:
        \begin{equation*}
            \sigma=\dot{e}+m\,e,\qquad m>0
        \end{equation*}

        Si $\sigma=0$, entonces $\dot{e}+m e=0$. Se propone la función candidata de Lyapunov:

        \begin{equation*}
            V=\frac{1}{2}\sigma^2
        \end{equation*}

        Derivando y sustituyendo las expresiones del error:

        \begin{equation*}
            \begin{aligned}
                \dot{\sigma} &= \ddot{e}+m\dot{e} \\
                &= \ddot{r}-f(X,t)-v+m(\dot{r}-X_2) \\
                &= m\dot{r}+\ddot{r}-mX_2-f(X,t)-v
            \end{aligned}
        \end{equation*}

        Considerando que las derivadas de la referencia y la dinámica no lineal están acotadas:
        \begin{equation*}
            |m\dot{r}+\ddot{r}|\le R,\qquad |f(X,t)|\le L
        \end{equation*}
        se tiene que:
        \begin{equation*}
            \dot{V}=\sigma\dot{\sigma}\le\sigma\left[R+L-mX_2-v\right]
        \end{equation*}

        Se elige la ley de control:
        \begin{equation*}
            v=-mX_2+K\operatorname{sign}(\sigma),\qquad K>0
        \end{equation*}

        Sustituyendo:
        \begin{equation*}
            \begin{aligned}
                \dot{V} &\le \sigma\left[R+L-mX_2+mX_2-K\operatorname{sign}(\sigma)\right] \\
                &\le (R+L)|\sigma|-K|\sigma| \\
                &\le -(K-R-L)|\sigma|
            \end{aligned}
        \end{equation*}

        Si $K>R+L$, entonces $\dot{V}<0$ para todo $\sigma\neq 0$, y el error de seguimiento es asintóticamente estable, es decir, $e\to 0$ conforme $t\to\infty$. Aplicando el resultado a cada eje, se obtiene:

        \begin{equation*}
            \begin{aligned}
                \sigma_x &= \dot{e}_x+n\,e_x, & v_x &= -n\dot{x}+l\operatorname{sign}(\sigma_x) \\
                \sigma_z &= \dot{e}_z+m\,e_z, & v_z &= -m\dot{z}+k\operatorname{sign}(\sigma_z)
            \end{aligned}
        \end{equation*}
        con $n,m>0$ y $l,k>R+L$.

        La dinámica de la actitud es:
        \begin{equation*}
            \ddot{\theta}=\frac{1}{J}\tau
        \end{equation*}
        donde $\theta_d$ es la referencia generada por el lazo externo. Se define el error de actitud:
        \begin{equation*}
            e_\theta = \theta-\theta_d
        \end{equation*}
        y la superficie de deslizamiento:
        \begin{equation*}
            \sigma_\theta = \dot{e}_\theta + b\,e_\theta,\qquad b>0
        \end{equation*}

        Se propone la función candidata de Lyapunov:
        \begin{equation*}
            V_\theta=\frac{1}{2}\sigma_\theta^2
        \end{equation*}
        cuya derivada es:
        \begin{equation*}
            \dot{V}_\theta = \sigma_\theta\dot{\sigma}_\theta
        \end{equation*}
        Derivando la superficie:
        \begin{equation*}
            \dot{\sigma}_\theta = \ddot{e}_\theta+b\dot{e}_\theta = \ddot{\theta}-\ddot{\theta}_d+b(\dot{\theta}-\dot{\theta}_d)
        \end{equation*}
        Sustituyendo la dinámica de actitud y eligiendo la ley de control:
        \begin{equation*}
            \tau = -b\dot{\theta}-j\operatorname{sign}(\sigma_\theta),\qquad j>0
        \end{equation*}
        se obtiene:
        \begin{equation*}
            \begin{aligned}
                \dot{\sigma}_\theta &= \frac{1}{J}\left(-b\dot{\theta}-j\operatorname{sign}(\sigma_\theta)\right)-\ddot{\theta}_d+b\dot{\theta}-b\dot{\theta}_d \\
                &= -\frac{j}{J}\operatorname{sign}(\sigma_\theta)-\left(\ddot{\theta}_d+b\dot{\theta}_d\right)
            \end{aligned}
        \end{equation*}
        y por lo tanto:
        \begin{equation*}
            \dot{V}_\theta \le -\frac{j}{J}|\sigma_\theta|+\left|\ddot{\theta}_d+b\dot{\theta}_d\right||\sigma_\theta|
        \end{equation*}
        Considerando que $|\ddot{\theta}_d+b\dot{\theta}_d|\le L_\theta$, se tiene:
        \begin{equation*}
            \dot{V}_\theta \le -\frac{1}{J}\left(j-JL_\theta\right)|\sigma_\theta|
        \end{equation*}
        Si $j>J L_\theta$, entonces $\dot{V}_\theta<0$, la superficie $\sigma_\theta\to 0$ en tiempo finito, y en consecuencia $e_\theta\to 0$ y $\dot{e}_\theta\to 0$, es decir, $\theta\to\theta_d$ y $\dot{\theta}\to\dot{\theta}_d$.

        


    \end{itemize}
\end{enumerate}

\end{document}
